\chapter{Chinese Remainder Theory}
\label{ch:v04}

Chinese remainder theory decomposes rings modulo comaximal ideals into products. It is both a structural theorem and an explicit congruence-solving method.

\section*{Learning goals}
After this chapter, the reader should be able to:
\begin{itemize}
\item verify comaximality of ideals and apply the Chinese remainder theorem;
\item construct explicit isomorphisms between quotient rings and products of quotients;
\item solve simultaneous congruences by lifting idempotents or Bezout relations;
\item extend the two-ideal theorem to finite pairwise-comaximal families;
\item compute kernels of diagonal quotient maps and identify intersections of ideals;
\item distinguish direct-product decompositions that arise from comaximal ideals from unrelated product presentations;
\end{itemize}
\section*{Conceptual roadmap}
\[
\boxed{\text{presentation or universal property}\;\longrightarrow\;\text{exactness/localization}\;\longrightarrow\;\text{structural consequence}.}
\]

\section{Comaximal ideals}
Ideals $I,J$ are comaximal when $I+J=A$.

\section{Two-ideal CRT}
For comaximal $I,J$, the quotient by $I\cap J$ maps isomorphically to $A/I\times A/J$.

% BEGIN VOL05-EXPANSION V04-example-01
\begin{example}[Integer CRT]\label{ex:v04-expand-01}
The congruences $x\equiv2\pmod3$ and $x\equiv4\pmod5$ have the unique solution $x\equiv14\pmod{15}$. This is the concrete form of $\mathbb Z/15\cong\mathbb Z/3\times\mathbb Z/5$.
\end{example}
% END VOL05-EXPANSION V04-example-01
\section{Product equals intersection}
For comaximal ideals, $I\cap J=IJ$.

\section{Finite families}
Pairwise comaximal finite families give product decompositions modulo the intersection or product.

\section{Explicit idempotents}
CRT decompositions produce orthogonal idempotents selecting the different factors.

% BEGIN VOL05-EXPANSION V04-example-02
\begin{example}[Orthogonal idempotents]\label{ex:v04-expand-02}
In $\mathbb Z/12\cong\mathbb Z/3\times\mathbb Z/4$, the classes $4$ and $9$ correspond to $(1,0)$ and $(0,1)$. They are orthogonal idempotents and sum to $1$.
\end{example}
% END VOL05-EXPANSION V04-example-02
\section{Integer congruences}
The classical integer CRT is the special case of principal ideals in $\mathbb Z$.

\section{Polynomial interpolation form}
Comaximal polynomial ideals encode simultaneous residue conditions.

% BEGIN VOL05-EXPANSION V04-example-03
\begin{example}[Polynomial CRT]\label{ex:v04-expand-03}
Since $(x)+(x-1)=k[x]$, one has $k[x]/(x(x-1))\cong k\times k$. A polynomial class is determined by its values at $0$ and $1$.
\end{example}
% END VOL05-EXPANSION V04-example-03
\section{Decomposition philosophy}
Product rings are algebraic manifestations of independent components.

\section{Core structural results}

\begin{theorem}[Two-ideal Chinese remainder theorem]\label{thm:v04-01}
If $I+J=A$, then $A/(I\cap J)\cong A/I\times A/J$.
\begin{proof}
The natural map has kernel $I\cap J$. For surjectivity choose $u\in I,v\in J$ with $u+v=1$; the element $bv+au$ has prescribed residues $a$ modulo $I$ and $b$ modulo $J$.
\end{proof}
\end{theorem}

\begin{theorem}[Intersection-product equality]\label{thm:v04-02}
If $I+J=A$, then $I\cap J=IJ$.
\begin{proof}
The inclusion $IJ\subseteq I\cap J$ is automatic. If $x\in I\cap J$ and $1=u+v$ with $u\in I,v\in J$, then $x=xu+xv$ lies in $IJ$.
\end{proof}
\end{theorem}

\begin{theorem}[Finite CRT]\label{thm:v04-03}
If $I_1,\dots,I_n$ are pairwise comaximal, then $A/(I_1\cdots I_n)\cong\prod_i A/I_i$.
\begin{proof}
Induct on $n$, using that the product of the first $n-1$ ideals is comaximal with $I_n$, then apply the two-ideal theorem and the product-intersection identity.
\end{proof}
\end{theorem}

\section{Worked examples}

\begin{example}[Comaximal ideals diagnostic]\label{ex:v04-worked-01}
Give a rigorous worked analysis of Comaximal ideals. State the ring, module, finiteness, localization, or homological hypotheses and derive the central conclusion. Ideals $I,J$ are comaximal when $I+J=A$. The decisive step is to use the universal property, exactness criterion, localization test, or finite-generation mechanism appropriate to the algebraic structure.
\end{example}

\begin{example}[Two-ideal CRT diagnostic]\label{ex:v04-worked-02}
Give a rigorous worked analysis of Two-ideal CRT. State the ring, module, finiteness, localization, or homological hypotheses and derive the central conclusion. For comaximal $I,J$, the quotient by $I\cap J$ maps isomorphically to $A/I\times A/J$. The decisive step is to use the universal property, exactness criterion, localization test, or finite-generation mechanism appropriate to the algebraic structure.
\end{example}

\begin{example}[Product equals intersection diagnostic]\label{ex:v04-worked-03}
Give a rigorous worked analysis of Product equals intersection. State the ring, module, finiteness, localization, or homological hypotheses and derive the central conclusion. For comaximal ideals, $I\cap J=IJ$. The decisive step is to use the universal property, exactness criterion, localization test, or finite-generation mechanism appropriate to the algebraic structure.
\end{example}

\begin{example}[Finite families diagnostic]\label{ex:v04-worked-04}
Give a rigorous worked analysis of Finite families. State the ring, module, finiteness, localization, or homological hypotheses and derive the central conclusion. Pairwise comaximal finite families give product decompositions modulo the intersection or product. The decisive step is to use the universal property, exactness criterion, localization test, or finite-generation mechanism appropriate to the algebraic structure.
\end{example}

\section{Solved dossiers}
Each dossier is a canonical solved problem. The provenance ledger records which retained corpus rules guided its topic and which dossiers were added freshly to complete the chapter.

\begin{problem}[Comaximal ideals diagnostic]\label{prob:v04-dossier-01}
Give a rigorous worked analysis of Comaximal ideals. State the ring, module, finiteness, localization, or homological hypotheses and derive the central conclusion.
\end{problem}
\begin{solution}
Ideals $I,J$ are comaximal when $I+J=A$. The decisive step is to use the universal property, exactness criterion, localization test, or finite-generation mechanism appropriate to the algebraic structure.
\end{solution}

\begin{problem}[Two-ideal CRT diagnostic]\label{prob:v04-dossier-02}
Give a rigorous worked analysis of Two-ideal CRT. State the ring, module, finiteness, localization, or homological hypotheses and derive the central conclusion.
\end{problem}
\begin{solution}
For comaximal $I,J$, the quotient by $I\cap J$ maps isomorphically to $A/I\times A/J$. The decisive step is to use the universal property, exactness criterion, localization test, or finite-generation mechanism appropriate to the algebraic structure.
\end{solution}

\begin{problem}[Product equals intersection diagnostic]\label{prob:v04-dossier-03}
Give a rigorous worked analysis of Product equals intersection. State the ring, module, finiteness, localization, or homological hypotheses and derive the central conclusion.
\end{problem}
\begin{solution}
For comaximal ideals, $I\cap J=IJ$. The decisive step is to use the universal property, exactness criterion, localization test, or finite-generation mechanism appropriate to the algebraic structure.
\end{solution}

\begin{problem}[Finite families diagnostic]\label{prob:v04-dossier-04}
Give a rigorous worked analysis of Finite families. State the ring, module, finiteness, localization, or homological hypotheses and derive the central conclusion.
\end{problem}
\begin{solution}
Pairwise comaximal finite families give product decompositions modulo the intersection or product. The decisive step is to use the universal property, exactness criterion, localization test, or finite-generation mechanism appropriate to the algebraic structure.
\end{solution}

\begin{problem}[Explicit idempotents diagnostic]\label{prob:v04-dossier-05}
Give a rigorous worked analysis of Explicit idempotents. State the ring, module, finiteness, localization, or homological hypotheses and derive the central conclusion.
\end{problem}
\begin{solution}
CRT decompositions produce orthogonal idempotents selecting the different factors. The decisive step is to use the universal property, exactness criterion, localization test, or finite-generation mechanism appropriate to the algebraic structure.
\end{solution}

\begin{problem}[Integer congruences diagnostic]\label{prob:v04-dossier-06}
Give a rigorous worked analysis of Integer congruences. State the ring, module, finiteness, localization, or homological hypotheses and derive the central conclusion.
\end{problem}
\begin{solution}
The classical integer CRT is the special case of principal ideals in $\mathbb Z$. The decisive step is to use the universal property, exactness criterion, localization test, or finite-generation mechanism appropriate to the algebraic structure.
\end{solution}

\begin{problem}[Polynomial interpolation form diagnostic]\label{prob:v04-dossier-07}
Give a rigorous worked analysis of Polynomial interpolation form. State the ring, module, finiteness, localization, or homological hypotheses and derive the central conclusion.
\end{problem}
\begin{solution}
Comaximal polynomial ideals encode simultaneous residue conditions. The decisive step is to use the universal property, exactness criterion, localization test, or finite-generation mechanism appropriate to the algebraic structure.
\end{solution}

\begin{problem}[Decomposition philosophy diagnostic]\label{prob:v04-dossier-08}
Give a rigorous worked analysis of Decomposition philosophy. State the ring, module, finiteness, localization, or homological hypotheses and derive the central conclusion.
\end{problem}
\begin{solution}
Product rings are algebraic manifestations of independent components. The decisive step is to use the universal property, exactness criterion, localization test, or finite-generation mechanism appropriate to the algebraic structure.
\end{solution}

\begin{problem}[Two-ideal Chinese remainder theorem proof dossier]\label{prob:v04-dossier-09}
Prove the following structural statement and identify the hypothesis that makes the conclusion work: If $I+J=A$, then $A/(I\cap J)\cong A/I\times A/J$.
\end{problem}
\begin{solution}
The natural map has kernel $I\cap J$. For surjectivity choose $u\in I,v\in J$ with $u+v=1$; the element $bv+au$ has prescribed residues $a$ modulo $I$ and $b$ modulo $J$. This proof isolates the reusable algebraic mechanism behind the result.
\end{solution}

\begin{problem}[Intersection-product equality proof dossier]\label{prob:v04-dossier-10}
Prove the following structural statement and identify the hypothesis that makes the conclusion work: If $I+J=A$, then $I\cap J=IJ$.
\end{problem}
\begin{solution}
The inclusion $IJ\subseteq I\cap J$ is automatic. If $x\in I\cap J$ and $1=u+v$ with $u\in I,v\in J$, then $x=xu+xv$ lies in $IJ$. This proof isolates the reusable algebraic mechanism behind the result.
\end{solution}

\begin{problem}[Finite CRT proof dossier]\label{prob:v04-dossier-11}
Prove the following structural statement and identify the hypothesis that makes the conclusion work: If $I_1,\dots,I_n$ are pairwise comaximal, then $A/(I_1\cdots I_n)\cong\prod_i A/I_i$.
\end{problem}
\begin{solution}
Induct on $n$, using that the product of the first $n-1$ ideals is comaximal with $I_n$, then apply the two-ideal theorem and the product-intersection identity. This proof isolates the reusable algebraic mechanism behind the result.
\end{solution}

\begin{problem}[Chapter synthesis]\label{prob:v04-dossier-12}
Connect the definitions and principal structural theorems of this chapter into one reusable commutative-algebra strategy.
\end{problem}
\begin{solution}
Chinese remainder theory decomposes rings modulo comaximal ideals into products. It is both a structural theorem and an explicit congruence-solving method. The recurring strategy is to encode the algebraic object by kernels, quotients, localization, finite presentation, or a resolution, apply the corresponding universal or exactness principle, and then recover the desired global or local conclusion.
\end{solution}

\section{Exercises with complete solutions}

\begin{exercise}\label{exr:v04-01}
State and apply the central principle behind Comaximal ideals. Give one concrete algebraic consequence.
\end{exercise}
\begin{hint}
Prove comaximality by finding a+b=1 with a in I and b in J; this identity drives the explicit CRT decomposition.
\end{hint}
\begin{solution}
Ideals $I,J$ are comaximal when $I+J=A$. The same principle can then be reused in later localization, support, base-change, resolution, Tor, or Ext arguments.
\end{solution}

\begin{exercise}\label{exr:v04-02}
State and apply the central principle behind Two-ideal CRT. Give one concrete algebraic consequence.
\end{exercise}
\begin{hint}
The kernel of A -> A/I x A/J is I intersect J; compute it before invoking the first isomorphism theorem.
\end{hint}
\begin{solution}
For comaximal $I,J$, the quotient by $I\cap J$ maps isomorphically to $A/I\times A/J$. The same principle can then be reused in later localization, support, base-change, resolution, Tor, or Ext arguments.
\end{solution}

\begin{exercise}\label{exr:v04-03}
State and apply the central principle behind Product equals intersection. Give one concrete algebraic consequence.
\end{exercise}
\begin{hint}
When I+J=A, prove I intersect J=IJ by inserting 1=a+b into an arbitrary element of the intersection.
\end{hint}
\begin{solution}
For comaximal ideals, $I\cap J=IJ$. The same principle can then be reused in later localization, support, base-change, resolution, Tor, or Ext arguments.
\end{solution}

\begin{exercise}\label{exr:v04-04}
State and apply the central principle behind Finite families. Give one concrete algebraic consequence.
\end{exercise}
\begin{hint}
To construct a simultaneous solution, build idempotent-like coefficients from a Bezout relation and combine the prescribed residues.
\end{hint}
\begin{solution}
Pairwise comaximal finite families give product decompositions modulo the intersection or product. The same principle can then be reused in later localization, support, base-change, resolution, Tor, or Ext arguments.
\end{solution}

\begin{exercise}\label{exr:v04-05}
State and apply the central principle behind Explicit idempotents. Give one concrete algebraic consequence.
\end{exercise}
\begin{hint}
For finitely many ideals, verify pairwise comaximality and use induction rather than assuming the two-ideal proof automatically scales.
\end{hint}
\begin{solution}
CRT decompositions produce orthogonal idempotents selecting the different factors. The same principle can then be reused in later localization, support, base-change, resolution, Tor, or Ext arguments.
\end{solution}

\begin{exercise}\label{exr:v04-06}
State and apply the central principle behind Integer congruences. Give one concrete algebraic consequence.
\end{exercise}
\begin{hint}
In Z, translate comaximality to gcd(m,n)=1 before solving congruences.
\end{hint}
\begin{solution}
The classical integer CRT is the special case of principal ideals in $\mathbb Z$. The same principle can then be reused in later localization, support, base-change, resolution, Tor, or Ext arguments.
\end{solution}

\begin{exercise}\label{exr:v04-07}
State and apply the central principle behind Polynomial interpolation form. Give one concrete algebraic consequence.
\end{exercise}
\begin{hint}
Check surjectivity of the product quotient map by lifting one coordinate at a time with the CRT coefficients.
\end{hint}
\begin{solution}
Comaximal polynomial ideals encode simultaneous residue conditions. The same principle can then be reused in later localization, support, base-change, resolution, Tor, or Ext arguments.
\end{solution}

\begin{exercise}\label{exr:v04-08}
State and apply the central principle behind Decomposition philosophy. Give one concrete algebraic consequence.
\end{exercise}
\begin{hint}
Do not replace intersections by products without comaximality; test a simple counterexample such as powers of one prime ideal.
\end{hint}
\begin{solution}
Product rings are algebraic manifestations of independent components. The same principle can then be reused in later localization, support, base-change, resolution, Tor, or Ext arguments.
\end{solution}

\section*{Chapter summary}
The chapter now contributes a canonical solved layer to the progression from rings and modules through localization, Noetherian methods, integral dependence, and derived functors.
% BEGIN VOL05-EXPANSION V04-exercises-01
\section{Graded supplementary exercises}
These exercises extend the protected chapter with a balanced set of computations, proofs, hypothesis tests, applications, and challenges.

\subsection*{Standard computations}

\begin{exercise}[Two congruences]\label{exr:v04-expand-01}
Solve $x\equiv1\pmod4$ and $x\equiv2\pmod3$.
\end{exercise}
\begin{hint}
Use comaximal ideals and Chinese remainder decompositions.
\end{hint}
\begin{solution}
$x\equiv5\pmod{12}$.
\end{solution}

\begin{exercise}[CRT decomposition]\label{exr:v04-expand-02}
Decompose $\mathbb Z/35\mathbb Z$.
\end{exercise}
\begin{hint}
Use comaximal ideals and Chinese remainder decompositions.
\end{hint}
\begin{solution}
$\mathbb Z/35\cong\mathbb Z/5\times\mathbb Z/7$.
\end{solution}

\begin{exercise}[Comaximality]\label{exr:v04-expand-03}
Are $(x)$ and $(x-2)$ comaximal when $2$ is invertible?
\end{exercise}
\begin{hint}
Use comaximal ideals and Chinese remainder decompositions.
\end{hint}
\begin{solution}
Yes.
\end{solution}

\begin{exercise}[Intersection]\label{exr:v04-expand-04}
Compute $(3)\cap(5)$ in $\mathbb Z$.
\end{exercise}
\begin{hint}
Use comaximal ideals and Chinese remainder decompositions.
\end{hint}
\begin{solution}
It is $(15)$.
\end{solution}

\begin{exercise}[Idempotent]\label{exr:v04-expand-05}
Find the class mod $15$ corresponding to $(1,0)$ mod $(3,5)$.
\end{exercise}
\begin{hint}
Use comaximal ideals and Chinese remainder decompositions.
\end{hint}
\begin{solution}
The class $10$.
\end{solution}

\subsection*{Proofs}

\begin{exercise}[Two-ideal Chinese remainder theorem proof]\label{exr:v04-expand-06}
Prove: If $I+J=A$, then $A/(I\cap J)\cong A/I\times A/J$.
\end{exercise}
\begin{hint}
Start from the chapter theorem hypotheses and identify the decisive algebraic map or containment.
\end{hint}
\begin{solution}
The natural map has kernel $I\cap J$. For surjectivity choose $u\in I,v\in J$ with $u+v=1$; the element $bv+au$ has prescribed residues $a$ modulo $I$ and $b$ modulo $J$.
\end{solution}

\begin{exercise}[Intersection-product equality proof]\label{exr:v04-expand-07}
Prove: If $I+J=A$, then $I\cap J=IJ$.
\end{exercise}
\begin{hint}
Start from the chapter theorem hypotheses and identify the decisive algebraic map or containment.
\end{hint}
\begin{solution}
The inclusion $IJ\subseteq I\cap J$ is automatic. If $x\in I\cap J$ and $1=u+v$ with $u\in I,v\in J$, then $x=xu+xv$ lies in $IJ$.
\end{solution}

\begin{exercise}[Finite CRT proof]\label{exr:v04-expand-08}
Prove: If $I_1,\dots,I_n$ are pairwise comaximal, then $A/(I_1\cdots I_n)\cong\prod_i A/I_i$.
\end{exercise}
\begin{hint}
Start from the chapter theorem hypotheses and identify the decisive algebraic map or containment.
\end{hint}
\begin{solution}
Induct on $n$, using that the product of the first $n-1$ ideals is comaximal with $I_n$, then apply the two-ideal theorem and the product-intersection identity.
\end{solution}

\begin{exercise}[Structural synthesis]\label{exr:v04-expand-09}
Prove a short corollary by combining Two-ideal Chinese remainder theorem with Intersection-product equality. State every hypothesis you use.
\end{exercise}
\begin{hint}
Apply the first theorem to move to its structural description, then use the second theorem on that description.
\end{hint}
\begin{solution}
A valid synthesis first invokes Two-ideal Chinese remainder theorem and then applies Intersection-product equality; the conclusion follows after checking the shared hypotheses.
\end{solution}

\subsection*{Counterexamples and hypothesis tests}

\begin{exercise}[Noncomaximal CRT]\label{exr:v04-expand-10}
Test the claim: $\mathbb Z/8\cong\mathbb Z/4\times\mathbb Z/2$.
\end{exercise}
\begin{hint}
Try a smallest standard ring or module from this chapter.
\end{hint}
\begin{solution}
False: the ideals are not comaximal.
\end{solution}

\begin{exercise}[Intersection product]\label{exr:v04-expand-11}
Test the claim: $I\cap J=IJ$ for all ideals.
\end{exercise}
\begin{hint}
Try a smallest standard ring or module from this chapter.
\end{hint}
\begin{solution}
False: take $I=J=(2)$ in $\mathbb Z$.
\end{solution}

\begin{exercise}[Repeated factors]\label{exr:v04-expand-12}
Test the claim: CRT works unchanged with repeated equal proper ideals.
\end{exercise}
\begin{hint}
Try a smallest standard ring or module from this chapter.
\end{hint}
\begin{solution}
False: the diagonal map is not surjective.
\end{solution}

\subsection*{Applications and investigations}

\begin{exercise}[Parallel arithmetic]\label{exr:v04-expand-13}
Why can CRT split modular arithmetic into components?
\end{exercise}
\begin{hint}
Translate the situation into comaximal ideals and Chinese remainder decompositions.
\end{hint}
\begin{solution}
The product-ring isomorphism makes operations componentwise.
\end{solution}

\begin{exercise}[Interpolation]\label{exr:v04-expand-14}
Interpret CRT for distinct linear polynomial factors.
\end{exercise}
\begin{hint}
Translate the situation into comaximal ideals and Chinese remainder decompositions.
\end{hint}
\begin{solution}
It prescribes independent values at distinct points.
\end{solution}

\subsection*{Challenge problems}

\begin{exercise}[Local-global synthesis]\label{exr:v04-expand-15}
Connect the chapter ideas 'Comaximal ideals' and 'Decomposition philosophy' in one rigorous argument.
\end{exercise}
\begin{hint}
State the relevant map, ideal, module, or universal property before drawing the conclusion.
\end{hint}
\begin{solution}
The argument begins with Ideals $I,J$ are comaximal when $I+J=A$. It then uses the later viewpoint: Product rings are algebraic manifestations of independent components. The bridge is supplied by the structural theorems proved in the chapter.
\end{solution}

\begin{exercise}[Hypothesis audit]\label{exr:v04-expand-16}
Choose one theorem from the chapter and explain precisely what can fail if one key hypothesis is removed.
\end{exercise}
\begin{hint}
Use one of the explicit counterexamples in the graded set and compare it with the theorem statement.
\end{hint}
\begin{solution}
The theorem hypotheses are essential because the chapter's quotient, localization, exactness, or finiteness mechanism can fail outside them. A correct answer identifies the dropped hypothesis and exhibits a concrete failure.
\end{solution}

% END VOL05-EXPANSION V04-exercises-01
